\documentclass[12pt,a4paper]{article}
\usepackage[utf8]{inputenc}
\usepackage[spanish,german]{babel}
\usepackage[T1]{fontenc}
\usepackage{geometry}
\geometry{a4paper, margin=2.5cm}
\usepackage{titlesec}
\usepackage{enumitem}
\usepackage{multicol}
\usepackage{csquotes}

\titleformat{\section}
{\Large\bfseries}
{}
{0em}
{}[\titlerule]

\titleformat{\subsection}
{\large\bfseries}
{}
{0em}
{}

\setlength{\parindent}{0pt}
\setlength{\parskip}{0.5em}

\begin{document}

\title{\textbf{DIALOG: Describir el Lugar donde Vives}}
\author{Spanisch A1/A2}
\date{}
\maketitle

\textbf{Tema:} Describir el lugar donde vives \\
\textbf{Nivel:} A1-A2 \\
\textbf{Complejidad:} Leicht-Mittel \\
\textbf{Palabras:} 892

\vspace{1em}

\textit{Nota: Se recomienda revisar primero el vocabulario antes de leer el diálogo.}

\vspace{1em}

\section*{Diálogo}

\textbf{Andrés:} Hola, Sofía. ¿Cómo estás?

\textbf{Sofía:} Hola, Andrés. Estoy muy bien, gracias. ¿Y tú?

\textbf{Andrés:} Muy bien también. Oye, quería preguntarte, ¿dónde vives exactamente?

\textbf{Sofía:} Vivo en un barrio muy bonito en el centro de la ciudad. Es un lugar tranquilo y seguro.

\textbf{Andrés:} ¡Qué bien! ¿Cómo es tu apartamento?

\textbf{Sofía:} Mi apartamento es pequeño pero muy cómodo. Tiene dos habitaciones, una cocina, un baño y una sala de estar. Es perfecto para mí.

\textbf{Andrés:} Suena muy agradable. ¿Hay muchos vecinos en tu edificio?

\textbf{Sofía:} Sí, hay muchos vecinos. El edificio es grande, tiene cinco pisos. Hay aproximadamente veinte apartamentos.

\textbf{Andrés:} ¿Y hay algún parque cerca de tu casa?

\textbf{Sofía:} Sí, hay un parque muy grande a dos cuadras de mi casa. Me gusta mucho ir allí los fines de semana.

\textbf{Andrés:} Eso es genial. ¿Hay supermercados cerca?

\textbf{Sofía:} Sí, hay varios supermercados cerca. También hay una farmacia, un banco y algunos restaurantes. Es muy conveniente.

\textbf{Andrés:} Perfecto. ¿Cómo está el transporte público en tu zona?

\textbf{Sofía:} Está muy bien. Hay una estación de metro a cinco minutos caminando. También hay varias líneas de autobús.

\textbf{Andrés:} Eso es excelente. ¿Tu barrio es ruidoso o tranquilo?

\textbf{Sofía:} Es bastante tranquilo durante el día, pero por la noche a veces hay mucho ruido porque hay muchos bares y restaurantes.

\textbf{Andrés:} Entiendo. ¿Y cómo está el ambiente en tu barrio? ¿Es amigable?

\textbf{Sofía:} Sí, está muy amigable. Los vecinos son muy simpáticos. Hay que saludar a todos cuando entras al edificio.

\textbf{Andrés:} Eso es muy bonito. ¿Hay problemas de seguridad en tu zona?

\textbf{Sofía:} No, no hay problemas. Es una zona muy segura. Hay que tener cuidado como en cualquier lugar, pero en general es muy tranquilo.

\textbf{Andrés:} Me alegra escuchar eso. ¿Y cómo es el clima en tu ciudad?

\textbf{Sofía:} El clima es muy agradable. En verano hace calor pero no demasiado. En invierno hace fresco pero no hace mucho frío. Es perfecto.

\textbf{Andrés:} Suena como un lugar ideal para vivir. ¿Cuánto tiempo llevas viviendo allí?

\textbf{Sofía:} Llevo viviendo allí dos años. Antes vivía con mis padres en otra parte de la ciudad.

\textbf{Andrés:} ¿Y estás contenta con tu nuevo lugar?

\textbf{Sofía:} Sí, estoy muy contenta. Me gusta mucho mi apartamento y mi barrio. Es un lugar perfecto para mí.

\textbf{Andrés:} Eso es genial. Oye, ¿hay algo que no te guste de tu barrio?

\textbf{Sofía:} Bueno, a veces hay que esperar mucho tiempo para el ascensor porque el edificio es alto. Pero en general, no hay nada que me moleste mucho.

\textbf{Andrés:} Eso no es tan malo. ¿Y hay espacios verdes cerca además del parque?

\textbf{Sofía:} Sí, hay algunos jardines pequeños en el barrio. También hay que caminar un poco para llegar a un parque más grande, pero está bien.

\textbf{Andrés:} Perfecto. Suena como un lugar muy agradable para vivir. Me gustaría visitarte algún día.

\textbf{Sofía:} ¡Por supuesto! Sería un placer. Mi apartamento está en la calle principal, es fácil de encontrar.

\textbf{Andrés:} Perfecto. Entonces nos vemos pronto. ¡Que tengas un buen día!

\textbf{Sofía:} Gracias, tú también. ¡Hasta pronto!

\newpage

\section*{Vocabulario}

\begin{multicols}{2}
\begin{itemize}[leftmargin=*]
    \item \textbf{barrio} - Stadtviertel, Viertel
    \item \textbf{tranquilo} - ruhig
    \item \textbf{seguro} - sicher
    \item \textbf{apartamento} - Wohnung, Apartment
    \item \textbf{cómodo} - bequem, komfortabel
    \item \textbf{habitación} - Zimmer
    \item \textbf{sala de estar} - Wohnzimmer
    \item \textbf{vecino} - Nachbar/Nachbarin
    \item \textbf{edificio} - Gebäude
    \item \textbf{piso} - Stockwerk, Etage
    \item \textbf{aproximadamente} - ungefähr
    \item \textbf{cuadra} - Block (Stadtblock)
    \item \textbf{conveniente} - praktisch, günstig
    \item \textbf{transporte público} - öffentlicher Verkehr
    \item \textbf{zona} - Zone, Gegend
    \item \textbf{estación de metro} - U-Bahn-Station
    \item \textbf{autobús} - Bus
    \item \textbf{ruidoso} - laut, geräuschvoll
    \item \textbf{ambiente} - Atmosphäre, Ambiente
    \item \textbf{amigable} - freundlich
    \item \textbf{simpático} - sympathisch, nett
    \item \textbf{saludar} - grüßen
    \item \textbf{seguridad} - Sicherheit
    \item \textbf{clima} - Klima, Wetter
    \item \textbf{agradable} - angenehm
    \item \textbf{llevar viviendo} - seit ... wohnen
    \item \textbf{contento} - zufrieden, glücklich
    \item \textbf{ascensor} - Aufzug
    \item \textbf{molestar} - stören, belästigen
    \item \textbf{espacios verdes} - Grünflächen
    \item \textbf{jardín} - Garten
    \item \textbf{calle principal} - Hauptstraße
\end{itemize}
\end{multicols}

\newpage

\section*{Preguntas de Comprensión}

\begin{enumerate}[leftmargin=*]
    \item ¿Dónde vive Sofía?
    \item ¿Cómo es el apartamento de Sofía?
    \item ¿Cuántas habitaciones tiene el apartamento?
    \item ¿Cuántos pisos tiene el edificio de Sofía?
    \item ¿Cuántos apartamentos hay aproximadamente en el edificio?
    \item ¿Dónde está el parque que menciona Sofía?
    \item ¿Qué hay cerca del apartamento de Sofía?
    \item ¿Dónde está la estación de metro?
    \item ¿Cómo es el barrio de Sofía durante el día?
    \item ¿Por qué a veces hay ruido por la noche?
    \item ¿Cómo son los vecinos según Sofía?
    \item ¿Hay problemas de seguridad en la zona de Sofía?
    \item ¿Cómo es el clima en la ciudad de Sofía?
    \item ¿Cuánto tiempo lleva Sofía viviendo en ese lugar?
    \item ¿Dónde vivía Sofía antes?
    \item ¿Está Sofía contenta con su nuevo lugar?
    \item ¿Qué no le gusta a Sofía de su barrio?
    \item ¿Qué hay además del parque cerca de su casa?
    \item ¿En qué calle está el apartamento de Sofía?
    \item ¿Qué quiere hacer Andrés en el futuro?
\end{enumerate}

\newpage

\section*{Verbo Irregular: TENER (haben)}

\textit{Nota: \enquote{vosotros/as} se usa solo en España, no en español sudamericano. Se incluye aquí para completitud.}

\begin{multicols}{2}

\subsection*{Presente}
\begin{itemize}[leftmargin=*]
    \item yo - tengo
    \item tú - tienes
    \item él/ella/usted - tiene
    \item nosotros/as - tenemos
    \item vosotros/as - tenéis
    \item ustedes - tienen
    \item ellos/ellas - tienen
\end{itemize}

\subsection*{Pretérito Indefinido}
\begin{itemize}[leftmargin=*]
    \item yo - tuve
    \item tú - tuviste
    \item él/ella/usted - tuvo
    \item nosotros/as - tuvimos
    \item vosotros/as - tuvisteis
    \item ustedes - tuvieron
    \item ellos/ellas - tuvieron
\end{itemize}

\subsection*{Pretérito Imperfecto}
\begin{itemize}[leftmargin=*]
    \item yo - tenía
    \item tú - tenías
    \item él/ella/usted - tenía
    \item nosotros/as - teníamos
    \item vosotros/as - teníais
    \item ustedes - tenían
    \item ellos/ellas - tenían
\end{itemize}

\columnbreak

\subsection*{Futuro Simple}
\begin{itemize}[leftmargin=*]
    \item yo - tendré
    \item tú - tendrás
    \item él/ella/usted - tendrá
    \item nosotros/as - tendremos
    \item vosotros/as - tendréis
    \item ustedes - tendrán
    \item ellos/ellas - tendrán
\end{itemize}

\subsection*{Pretérito Perfecto}
\begin{itemize}[leftmargin=*]
    \item yo - he tenido
    \item tú - has tenido
    \item él/ella/usted - ha tenido
    \item nosotros/as - hemos tenido
    \item vosotros/as - habéis tenido
    \item ustedes - han tenido
    \item ellos/ellas - han tenido
\end{itemize}

\end{multicols}

\newpage

\section*{Explicación Gramatical}

\textbf{Ser vs. Estar -- Der Unterschied zwischen \enquote{ser} und \enquote{estar}}

Beide Verben bedeuten \enquote{sein}, werden aber in unterschiedlichen Situationen verwendet. Der Unterschied ist fundamental für das Spanische.

\textbf{\enquote{SER} wird verwendet für:}
\begin{itemize}[leftmargin=*]
    \item \textbf{Eigenschaften und Charakteristika} (dauerhaft): \enquote{El apartamento es pequeño} (Die Wohnung ist klein)
    \item \textbf{Identität und Beruf}: \enquote{Soy estudiante} (Ich bin Student)
    \item \textbf{Herkunft und Nationalität}: \enquote{Soy de Colombia} (Ich bin aus Kolumbien)
    \item \textbf{Material}: \enquote{La mesa es de madera} (Der Tisch ist aus Holz)
    \item \textbf{Uhrzeit und Datum}: \enquote{Es la una} (Es ist ein Uhr)
    \item \textbf{Preis}: \enquote{Es caro} (Es ist teuer)
\end{itemize}

\textbf{\enquote{ESTAR} wird verwendet für:}
\begin{itemize}[leftmargin=*]
    \item \textbf{Zustand und Befindlichkeit} (temporär): \enquote{Estoy contenta} (Ich bin zufrieden)
    \item \textbf{Position und Ort}: \enquote{El parque está cerca} (Der Park ist in der Nähe)
    \item \textbf{Handlungen im Verlauf} (mit Gerundio): \enquote{Estoy estudiando} (Ich studiere gerade)
    \item \textbf{Temporäre Eigenschaften}: \enquote{Está enfermo} (Er ist krank)
\end{itemize}

\textbf{Beispiele aus dem Text:}
\begin{itemize}[leftmargin=*]
    \item \enquote{Mi apartamento es pequeño} -- Eigenschaft (SER)
    \item \enquote{Es un lugar tranquilo} -- Charakteristik (SER)
    \item \enquote{El edificio es grande} -- Eigenschaft (SER)
    \item \enquote{Estoy muy bien} -- Befindlichkeit (ESTAR)
    \item \enquote{Está muy bien} -- Zustand (ESTAR, bezogen auf Transport)
    \item \enquote{Está muy amigable} -- Zustand (ESTAR, bezogen auf Ambiente)
    \item \enquote{Estoy muy contenta} -- Befindlichkeit (ESTAR)
    \item \enquote{El apartamento está en la calle principal} -- Position (ESTAR)
\end{itemize}

\textbf{Die Struktur \enquote{hay} und \enquote{hay que}}

\textbf{\enquote{HAY} (es gibt, es sind):}
\begin{itemize}[leftmargin=*]
    \item Unpersönliche Form von \enquote{haber}
    \item Wird verwendet, um Existenz oder Vorhandensein auszudrücken
    \item Immer in der 3. Person Singular, auch bei Plural-Subjekten
    \item \enquote{Hay un parque} (Es gibt einen Park) / \enquote{Hay muchos vecinos} (Es gibt viele Nachbarn)
\end{itemize}

\textbf{\enquote{HAY QUE} (man muss, es ist nötig):}
\begin{itemize}[leftmargin=*]
    \item Drückt Notwendigkeit oder Verpflichtung aus
    \item Wird immer mit Infinitiv verwendet
    \item Unpersönlich -- kein Subjekt nötig
    \item \enquote{Hay que estudiar} (Man muss studieren)
\end{itemize}

\textbf{Beispiele aus dem Text:}
\begin{itemize}[leftmargin=*]
    \item \enquote{¿Hay muchos vecinos?} -- Gibt es viele Nachbarn?
    \item \enquote{Hay aproximadamente veinte apartamentos} -- Es gibt ungefähr zwanzig Wohnungen
    \item \enquote{¿Hay algún parque cerca?} -- Gibt es einen Park in der Nähe?
    \item \enquote{Hay un parque muy grande} -- Es gibt einen sehr großen Park
    \item \enquote{Hay varios supermercados} -- Es gibt mehrere Supermärkte
    \item \enquote{Hay una estación de metro} -- Es gibt eine U-Bahn-Station
    \item \enquote{Hay que saludar a todos} -- Man muss alle grüßen
    \item \enquote{Hay que tener cuidado} -- Man muss vorsichtig sein
    \item \enquote{Hay que esperar mucho tiempo} -- Man muss lange warten
    \item \enquote{Hay que caminar un poco} -- Man muss ein bisschen laufen
\end{itemize}

\textbf{Tipp}: \enquote{Ser} = dauerhaft, \enquote{Estar} = temporär oder Position. \enquote{Hay} = Existenz, \enquote{Hay que} = Notwendigkeit.

\end{document}
